%% ============================================================
%% Chapter 22: Sheaves of Necessity and Affordability
%% ============================================================

\chapter{Sheaves of Necessity and Affordability}
\label{ch:ch22-sheaves-of-necessity}
\section{What We Call Affordable}

A society reveals its actual moral ordering not in its stated values but in what it is willing to call affordable without argument, and in what it quietly decides must never become cheaper. Fireworks are affordable. A second school counselor is not. A stadium subsidy is affordable. A universal mentorship program is not. An industrial meat supply chain is affordable, in the sense that its true cost is never presented as a bill anyone has to approve, while a public health intervention with a fraction of that cost must survive a legislative session to be funded at all. And the falling prices that farming, manufacturing, and logistics improvements might otherwise have delivered are set aside instead, so that the real burden of existing debt can be eroded and asset valuations supported, a trade made permanently and by default rather than case by case in response to actual deflationary risk.

This is not an argument for joylessness, nor for monetary chaos, nor a demand that ordinary people give up rest, travel, the company of a pet, or a stable currency while distant poverty persists in some abstract, unreachable form. It is an argument that the discretionary and the necessary, and the stable and the exploitative, have been sorted by institutional habit rather than by honest accounting, and that habit has consistently favored spectacle, appetite, convenience, and the protection of existing claims over development and the diffusion of technological gain, because those are what the people who set budgets and write monetary policy already hold. The remedy is not universal austerity. It is the same accounting, applied consistently and without exception for familiarity: weigh external cost against available alternatives, apply that standard progressively so restraint falls where alternatives exist, and stop allowing the sheer familiarity of an expenditure or a policy default to substitute for its justification. Until that accounting is done, the words affordable and stable will keep doing quiet moral work that the society using them has never actually agreed to.

\section{Affordability as Admissibility}

What follows reinterprets the preceding argument in the vocabulary this monograph has used throughout for questions of constraint and admissibility. It introduces no claim the argument above could not already support on its own terms; it is a change of lens, applied at the close of the book to make explicit a structural pattern that has, in fact, already recurred at every scale examined so far.

Let $X$ be the space of possible allocations of a society's resources, and let $\mathcal{A}_t \subseteq X$ be the institutionally admissible region at time $t$: the set of allocations an institution will fund, permit, or subsidize without demanding fresh justification. Admissibility in this sense is a membership condition, not an optimality claim; an allocation can remain inside $\mathcal{A}_t$ indefinitely without ever having been evaluated against any openly stated welfare function, simply by virtue of having entered $\mathcal{A}_t$ historically and never having been reconsidered for removal. Stadium subsidies, industrial meat production, and spectacle infrastructure persist inside $\mathcal{A}_t$ for exactly this reason: normalization placed them there, and nothing in the ordinary operation of a budget cycle asks them to justify their continued membership. Mentorship, tutoring, and preventive care can carry substantially greater developmental value while remaining outside $\mathcal{A}_t$, unable to enter without discharging a proof burden that the already-admitted claims inside $\mathcal{A}_t$ never face even once.

This is precisely the admissible-but-non-committable structure Chapter~\ref{ch:ch10-clio-architecture}'s CLIO architecture formalized at the level of an individual agentic system: $\mathrm{Adm}(r)$ never by itself supplies $W(r)$ or $\mathrm{Authorized}(r)$, and a proposal's mere membership in a permitted region is not a warrant for its continuation. What CLIO treats as a design requirement for one acting system, this chapter treats as a diagnosis of an entire society's budget cycle: an allocation's presence inside $\mathcal{A}_t$ is a historical fact about what was once admitted, not a live judgment that it remains warranted, and the two are routinely conflated precisely because nothing in ordinary institutional operation forces the distinction into view. Genuine political change, on this reading, is not optimization within a fixed $\mathcal{A}_t$, reshuffling a feasible set whose boundary is treated as given; it is revision of the boundary itself, an operation on the region rather than within it, a modification of the operative constraint regime rather than a search confined to whatever that regime currently permits.

\section{The Category of Claims}

A second structure captures something admissibility alone does not: not every excluded claim is excluded for the same reason, and not every admitted claim arrived by the same route. Let $\mathbf{Claims}$ be a category whose objects are socially recognized claims on resources, and whose morphisms are accepted justifications, the argumentative or institutional moves by which one claim is converted into another, more legible institutional form: civic pride justifies a stadium; consumer choice justifies industrial agriculture; economic stability justifies an inflation target; property rights and the goal of attracting foreign investment justify a speculative housing purchase. This is a genuine category rather than a mere directed graph of associations because these justifications compose: if civic identity justifies tourism expenditure and tourism expenditure in turn justifies stadium infrastructure, their composition supplies an accepted route directly from civic identity to stadium funding, without either intermediate step needing to be re-argued, and every claim carries an identity morphism representing its bare recognition without translation into any other form.

The category is not richly connected. Some objects sit at the center of a dense web of morphisms with short compositional paths to funded outcomes; others are nearly isolated, reachable only through long, unfamiliar, or nonexistent chains of composition. A sports subsidy maps, by well-worn institutional paths, into employment policy, regional development, civic identity, tourism promotion, and public cohesion, any one of which supplies a sufficient justification on its own. A child's need for sustained mentorship typically has no comparably worn morphism into a funded budget line, and a tenant's claim to remain housed in their own neighborhood fares no better: property rights supply the buyer with an immediate, well-recognized morphism into legitimate ownership, while years of residence, employment, and community membership supply the tenant with no comparably accepted arrow into any form of institutional protection at all; the claim exists as an object, but the category contains no accepted route, direct or composite, carrying it there.

Privilege, in this vocabulary, is morphic abundance rather than intrinsic urgency: not merely many outgoing arrows, but many short compositional paths into the budgetary objects an institution already funds. A privileged claim need not be more pressing or more valuable than its excluded counterpart; it inhabits a richer justificatory neighborhood, one that institutions already have the translation apparatus, and the compositional shortcuts, to recognize, in the language of economics, culture, security, or tradition, whichever happens to be at hand. Deprivation persists in part because urgent claims lack comparable translation paths, not because the claims themselves are weaker. This is the categorical form of an asymmetry this monograph has met before in a different register: Chapter~\ref{ch:ch20-rebel-without-a-cost}'s account of contribution showed that a documented act does not compose, by any warranted morphism, into recognition, reward, or authority, $\mathrm{Record}(x) \not\Rightarrow \cdots \not\Rightarrow \mathrm{Continue}(x)$, and that a system minting one convertible path across all six acts manufactures exactly the compositional shortcut this chapter's category of claims shows privileged allocations already enjoy by historical accident rather than by argument. The entitled claim does not have to argue its case because it already has an accepted morphism, or a short composite of them, waiting; the developmental claim has to build one from nothing, in whichever vocabulary a given funding cycle happens to honor that year.

\section{A Sheaf of Necessities}

The deepest structural feature is neither the shape of the admissible region nor the density of justificatory morphisms, but a failure of composition across scale. Let $\{U_i\}$ be a cover of the shared resource space $X$ by institutional or social contexts, cities, households, industries, national governments, and let $\mathcal{N}(U_i)$ denote the collection of allocations regarded as necessary within context $U_i$, so that a particular such judgment is a local section $s_i \in \mathcal{N}(U_i)$. A city can regard a stadium as economically necessary; a household can regard an annual flight as personally necessary; an industry can regard cheap meat as commercially necessary; a government can regard restrained social expenditure as fiscally necessary. Each of these local sections can be internally coherent, and they can appear compatible on the limited overlaps that ordinary institutional practice actually inspects, a city's accounting rarely checks itself against a household's, an industry's rarely checks itself against a government's long-run fiscal position. But apparent pairwise compatibility on the overlaps actually inspected is weaker than genuine compatibility on the full cover. Once ecological, intergenerational, and developmental effects are included in the restriction maps, some of these sections disagree on overlaps that ordinary accounting had simply omitted from consideration.

The absence of a global section $s \in \mathcal{N}(X)$ compatible with ecological limits and universal developmental provision is therefore not a mysterious failure of the sheaf axiom, which guarantees gluing precisely when compatibility genuinely holds on every overlap; it is evidence that the locally asserted necessities were never actually compatible over the full cover, only over the narrower set of overlaps institutions habitually check. What looks affordable at every overlap actually inspected can still be globally impossible; what appears optional at every overlap actually inspected can still be globally indispensable. Contemporary institutions are reasonably reliable at validating claims within a bounded jurisdiction, a household, a city, an industry; what they lack is a restriction map wide enough to catch the overlaps, ecological, intergenerational, planetary, that would reveal the incompatibility their narrower accounting conceals.

This is not a new construction in this monograph, and naming its recurrence is the point of placing this chapter last. Chapter~\ref{ch:ch08-tartan-tiling} formalized the identical obstruction in a physical register, TARTAN's tile-gluing defect $K_{ij}$, and read Buchanan's Kochen--Specker contextuality result as an independent instance of the same pattern: locally consistent measurement contexts that nonetheless admit no globally coherent assignment, the local-coherence-without-global-section failure stated there in the vocabulary of sheaf theory and quantum foundations. Chapter~\ref{ch:ch19-ledgers-without-value} formalized a close relative of it in an economic register, calling it degeneracy rather than obstruction: a speculative system whose admissible set $\mathrm{Adm}(\mathcal{L})$ is large and internally consistent at the level of the ledger, while providing no selection pressure toward the productive subset $\mathrm{Adm}^+(\mathcal{L})$ within it. All three, TARTAN's tiles, the ledger's admissible set, and this chapter's cover of necessities, share the same mathematical shape: local well-formedness is cheap and reliably achieved; the property that actually matters, global coherence, positive productivity, universal compatible necessity, is not automatic, and no amount of local diligence substitutes for checking it directly. Chapter~\ref{ch:ch01-latency-of-evidence}'s dependency chain, $C_t \Rightarrow A_t \Rightarrow L_t \Rightarrow X_t$ but not the reverse, and Chapter~\ref{ch:ch15-compression-fallacy}'s non-identification of $L_T$ and $D_T$ are the same warning stated at the level of a single trace or a single representation rather than a cover of contexts: an earlier, cheaper property holding gives no guarantee about a later, more demanding one, and the gap between them is exactly where a system's actual behavior, or a society's actual moral ordering, turns out to live.

Read together, the three sections of this chapter trace one failure at three depths. Admissibility explains why familiar pleasures are never asked to justify themselves twice, while developmental need is asked every time. The category of claims explains why some exclusions look like weakness of argument when they are really poverty of translation. The sheaf of necessities explains why a society composed entirely of locally reasonable judgments can still arrive, without any single decision that looks wrong from where it was made, at a global order that is not defensible at all.

\section{Conclusion: The Recurring Shape}

This monograph opened, in Chapter~\ref{ch:ch01-latency-of-evidence}, with a trace that can be true and fully evidenced long before any system is equipped to recognize it as such, a gap between what exists and what is admitted. It closes with the same gap, restated at the scale of an entire civilization's budget: allocations that have been admitted without ever having been checked, and necessities that have never been admitted despite deserving to be. Between these two points, the same shape recurred at every register this book examined: the distinction-holonomy apparatus of Chapter~\ref{ch:ch06-distinction-holonomy} showed that a system's momentary state never determines whether it will continue to behave the same way under transport, exactly as this chapter's admissible region $\mathcal{A}_t$ never determines whether an allocation remains warranted; the CLIO architecture of Chapter~\ref{ch:ch10-clio-architecture} showed that verification constrains continuation without authorizing commitment, exactly as this chapter's category of claims shows that a claim's mere existence never supplies the morphism that would let it compose into a funded outcome; and the sheaf-theoretic obstruction examined independently in physics, in cryptocurrency markets, and now in the allocation of a society's resources shows that local coherence, however carefully and repeatedly checked, is never a substitute for verifying the one property, global compatibility, productive selection, universal admissible necessity, that a system's continued legitimacy actually depends on.

None of this was planned as a single argument at the outset of the corpus this monograph draws together. That the same obstruction keeps reappearing, independently motivated each time by a different domain's own internal questions, is not proof that the pattern is fundamental in some deeper metaphysical sense. It is evidence of something more modest and, for a closing chapter, more useful: that the questions this book has been asking, what persists, what may be repaired, what a representation actually preserves, what a record warrants, and what a society is entitled to call affordable, are not five different questions after all, but one question, asked patiently, in whichever vocabulary the object at hand happened to require.
